\documentclass[book.tex]{subfiles}

\begin{document}

\chapter{Introduction to Quantum Chemistry Methods}
\label{chap:hartree_fock}
\noindent\rule{11cm}{0.4pt} \\
\textbf{Prerequisites:} \autoref{chap:molecular_hamiltonian},  \\
\textbf{Difficulty Level:} ** \\
\noindent\rule{11cm}{0.4pt}
\vspace{5mm}

Quantum chemists have built up a large family of methods to approximately solve Schrodinger's equation for large chemical systems.

\section{Hartree-Fock Methods}

A very simple ansatz for the joint wave function of a $N$ electrong system is given by the product of per-atom wave functions.

\begin{align}
\psi(r_1,\dotsc,r_N) &= \phi(r_1)\dotsc \phi(r_N)
\end{align}

This ansatz is too simple to work well in practice. A slightly more sophisticated ansatz uses the notion of Slater determinants. The Slater determinant is given by
\begin{align}
\psi(x_1, x_2) &= \frac{1}{\sqrt{2}} [ \chi_1(x_1)\chi_2(x_2) - \chi_1(x_2)\chi_2(x_1)]
\end{align}

The full determinant for $N$ electrons is given by
\begin{align}
\psi &= \frac{1}{\sqrt{N!}}\begin{vmatrix}
    \chi_1(x_1) & \dotsc & \chi_N(x_1) \\
    \vdots & \ddots & \vdots \\
    \chi_1(x_N) & \dotsc & \chi_N(x_N)
    \end{vmatrix}
\end{align}

The Hartree-Fock method takes the ansatz that the $N$-body wave function is given by a slater determinant. The method then performs a variational approximation to find the best function in this class.

\section{Coupled Cluster}

Coupled cluster methods are routinely viewed as the gold-standard for classical quantum chemistry calculations and are often used to generate training data to parameterize faster quantum chemistry techniques.

The wavefunction for the coupled-cluster method takes the exponential ansatz

\begin{align}
| \psi \rangle &= e^T |\phi_0\rangle
\end{align}
Here $\phi_0$ is a reference wave-function which is given by a Slater determinant of orbitals as in Hartree-Fock. $T$ is called the cluster operator and has the form
\begin{align}
T &= T_1 + T_2 + \dotsc
\end{align}
where $T_1$ is the operator of single excitations, $T_2$ the operator of double excitations and so on. These operators are formulated in terms of creation and annihilation operators.
\end{document}
